% Obsahový zdroj: import_hodiny/37_Vlastnosti_kombinacnich_cisel.tex a odpovídající PDF (obě strany).
% P1–P5: zdrojové Příklady 1–5; C1–C5: zdrojové Procvičování 1–5.
% Matematický obsah zachován; výklad, odvození a řešení pedagogicky rozvedeny.
% Upřesnění a opravy: viz KMB_29-41_Kontrola.md; původní import zůstává beze změny.
% Rozpis: 4_rocnik/tematicke_plany/Rozpis_vyuky_MAT_4_rocnik.tex, hodina 37.
% Plán: 4_rocnik/tematicke_plany/SPS.pdf, oddíly 2.1–2.3.
% Sazba: sablony/pripravy/TEMPLATE_Priprava.tex a společný priprava.sty.
% Vlastní doplnění: didaktický výklad, zdůvodnění, metodika; žádné nové číslované úlohy.
% Učebnice/09_Kombinatorika nemá obsahovou kapitolu; místní kombinatorické listy ani Realisticky.cz nejsou dostupné.
\documentclass[11pt,a4paper]{article}
% Spolecny styl pro ucitelske pripravy MAT (Work/MAT)
% Verze 1.1, 2026-10-23
\NeedsTeXFormat{LaTeX2e}
\ProvidesPackage{priprava}[2026/10/23 v1.1 Sablona ucitelskych priprav MAT]

% Zakladni typografie a cestina
\RequirePackage[T1]{fontenc}
\RequirePackage[utf8]{inputenc}
\RequirePackage[czech]{babel}
\RequirePackage{lmodern}
\RequirePackage{microtype}
\RequirePackage{amsmath,amssymb}
\RequirePackage{xcolor}
\RequirePackage{enumitem}
\RequirePackage[most]{tcolorbox}
\RequirePackage{tabularx,booktabs,array}
\RequirePackage{fancyhdr}
\RequirePackage{needspace}
\RequirePackage{geometry}
\RequirePackage[hidelinks]{hyperref}

\geometry{a4paper,top=17mm,bottom=18mm,left=23mm,right=23mm}
\definecolor{notegray}{gray}{0.39}
\setlength{\parindent}{0pt}
\setlength{\parskip}{2.5pt plus 1pt minus 0.5pt}
\setlist[itemize]{topsep=2pt,itemsep=1pt,parsep=0pt,leftmargin=1.45em}
\setlist[enumerate]{topsep=2pt,itemsep=2pt,parsep=0pt,leftmargin=1.65em}
\setcounter{secnumdepth}{0}
\renewcommand{\arraystretch}{1.12}

% Zahlavi: vlevo rocnik | globalni tema, vpravo cislo: nazev hodiny
% Zapati zustava zcela prazdne.
\pagestyle{fancy}
\fancyhf{}
\renewcommand{\headrulewidth}{0.35pt}
\renewcommand{\footrulewidth}{0pt}
\setlength{\headheight}{14pt}
\newcommand{\nastavPripravu}[4]{%
  \fancyhead[L]{\small #1\enspace\textbar\enspace #2}%
  \fancyhead[R]{\small #3: #4}%
}

% Minimalisticke cernobile boxy. Kazda definice zvlast.
\tcbset{commonlesson/.style={enhanced,breakable=false,colback=white,colframe=black,
  colbacktitle=white,coltitle=black,boxrule=0.45pt,
  arc=0.8mm,outer arc=0.8mm,left=3mm,right=3mm,top=1.6mm,bottom=1.6mm,
  toptitle=1mm,bottomtitle=0.7mm,fonttitle=\bfseries,
  attach boxed title to top left={xshift=2.5mm,yshift=-1.15mm},
  boxed title style={colback=white,colframe=white,boxrule=0pt,left=1mm,right=1mm}}}
\newtcolorbox{zadaniBox}[1][]{commonlesson,title={Zadání},#1}
\newtcolorbox{prikladBox}[1][]{commonlesson,title={Příklad},#1}
\newtcolorbox{definiceBox}[1][]{commonlesson,title={Definice},#1}

% Spolecny box pro zadani a jeho reseni, bez mezery mezi castmi.
% Pouziti:
%   \begin{prikladReseniBox}
%     Zadani ulohy.
%     \tcblower
%     \begin{align*} ... \end{align*}
%   \end{prikladReseniBox}
% Spodni cast ma jemny sedy podklad a delici linku; lze v ni sazet
% vice radku matematiky. Obe casti maji jen jeden spolecny ramecek.
\tcbset{solutionlesson/.style={
  commonlesson,
  enhanced jigsaw,
  breakable,
  colbacklower=black!3,
  segmentation style={solid,draw=black!25,line width=0.35pt},
  fontlower=\small\color{notegray},
  before lower={\noindent\textbf{Řešení:}\enspace},
  leftlower=3mm,rightlower=3mm
}}
\newtcolorbox{prikladReseniBox}[1][]{solutionlesson,title={Příklad},#1}
\newtcolorbox{zadaniReseniBox}[1][]{solutionlesson,title={Zadání},#1}

% Pomucky pro autora pripravy: sedive poznamky; \reseni pro starsi pripravy
\newcommand{\poznamka}[1]{\par\begingroup\small\color{notegray}\itshape
  \textnormal{\textbf{Poznámka pro učitele.}}\enspace #1\par\endgroup}
\newcommand{\reseni}[1]{\par\begingroup\small\color{notegray}
  \textbf{Očekávané řešení:}\enspace #1\par\endgroup}
\newcommand{\blok}[1]{\needspace{5\baselineskip}\subsection*{#1}}
\newcommand{\vyrok}[1]{\uv{#1}}

\nastavPripravu{4. ročník}{Kombinatorika}{37}{Vlastnosti kombinačních čísel}
\begin{document}


V hodině 36 jsme krátili faktoriály. Dnes budeme rovnosti nejprve vysvětlovat počítáním týchž objektů dvěma způsoby. Výpočet potom poslouží jako druhá kontrola.
\poznamka{Předpoklady: význam $\binom nk$, disjunktní případy a pravidlo součinu. Cíl: žák zdůvodní symetrii, Pascalův vztah i součet počtů všech podmnožin a pozná, kdy lze tyto vztahy použít. P5 a C5 jsou volitelné rozšíření s úplným řešením, nikoli další povinné vzorce.}

\blok{Výklad a odvození}

\textbf{Symetrie.} Pro celá $n\geq0$ a $0\leq k\leq n$ platí
\[
 \binom nk=\binom n{n-k}.
\]
Ke každé $k$-prvkové skupině přiřaď její doplněk v pevné $n$-prvkové množině. Doplněk má $n-k$ prvků a zpětným doplněním získáme původní skupinu. Jde o jednoznačné vzájemné přiřazení. Algebraicky se pouze vymění faktoriály ve jmenovateli:
\[
 \frac{n!}{k!(n-k)!}=\frac{n!}{(n-k)!k!}.
\]
Symetrie platí i pro krajní $k=0,n$; obě strany jsou jedna.

\textbf{Pascalův vztah.} Pro celá $n\geq2$ a $1\leq k\leq n-1$ platí
\[
 \binom nk=\binom{n-1}{k-1}+\binom{n-1}k.
\]
Zvolme konkrétní osobu P z celkem $n$ osob. V první skupině komisí P je: přidáme k ní $k-1$ z ostatních $n-1$ osob. Ve druhé skupině P není: vybíráme všech $k$ z ostatních. Případy se nepřekrývají a pokrývají všechny komise, takže jejich počty sčítáme. Podmínky zaručují, že obě čísla na pravé straně jsou definována. Krajní hodnoty řešíme samostatně jako $\binom n0=\binom nn=1$.

\textbf{Součet počtů podmnožin.} Pro celé $n\geq0$ platí
\[
 \sum_{k=0}^{n}\binom nk=2^n.
\]
Značka $\sum$ znamená sečíst členy pro všechna celá $k$ od $0$ do $n$. Na levé straně třídíme podmnožiny podle velikosti. Na pravé straně každému z $n$ prvků přiřazujeme jednu ze dvou odpovědí „vybrat/nevybrat“. Oba postupy počítají všechny podmnožiny včetně prázdné a úplné. Pro $n=0$ je jedinou podmnožinou prázdná množina a rovnost dává $1=2^0$.


\poznamka{Průběh: 0–11 min symetrie a P1; 11–25 min Pascalův vztah s P2–P3; 25–33 min součet a P4; 33–42 min C1–C4; 42–45 min vysvětlení jedné identity vlastními slovy. P5 a C5 ponech jako zásobu nebo domácí rozšíření; oba důkazy by plnou hodinu přeplnily.}

\blok{Společně řešené příklady}

% Zdroj P1: import_hodiny/37_Vlastnosti_kombinacnich_cisel.tex, Příklad 1.
\begin{prikladReseniBox}[title={Příklad P1}]
Vypočti \(\binom{10}{8}\).
\tcblower
Symetrií \(\binom{10}{8}=\binom{10}{2}=45\).
\end{prikladReseniBox}


% Zdroj P2: import_hodiny/37_Vlastnosti_kombinacnich_cisel.tex, Příklad 2.
\begin{prikladReseniBox}[title={Příklad P2}]
Vypočti \(\binom94+\binom93\) jediným kombinačním číslem.
\tcblower
Pascalův vztah s \(n=10,k=4\): \(\binom94+\binom93=\binom{10}4=210\). Přímá kontrola: \(126+84=210\).
\end{prikladReseniBox}


% Zdroj P3: import_hodiny/37_Vlastnosti_kombinacnich_cisel.tex, Příklad 3.
\begin{prikladReseniBox}[title={Příklad P3}]
Vysvětli rovnost \(\binom{10}4=\binom94+\binom93\) na výběru čtyřčlenné komise.
\tcblower
V desetihlavé skupině pevně určíme osobu P. Komise bez P: čtyři členové z devíti, tedy $\binom94=126$. Komise s P: doplníme jen tři z ostatních devíti, tedy $\binom93=84$. Každá čtyřčlenná komise je právě v jednom z těchto případů. Proto $\binom{10}4=126+84=210$.
\end{prikladReseniBox}


% Zdroj P4: import_hodiny/37_Vlastnosti_kombinacnich_cisel.tex, Příklad 4.
\begin{prikladReseniBox}[title={Příklad P4}]
Kolik podmnožin má pětiprvková množina?
\tcblower
U každého z pěti prvků se nezávisle rozhodneme, zda patří do podmnožiny; \(2^5=32\). Současně je to \(\binom50+\cdots+\binom55\).
\end{prikladReseniBox}


% Zdroj P5: import_hodiny/37_Vlastnosti_kombinacnich_cisel.tex, Příklad 5.
\begin{prikladReseniBox}[title={Příklad P5}]
Pro celá $n\geq1$ a $1\leq k\leq n$ vyjádři $\binom nk$ pomocí $\binom{n-1}{k-1}$ a vztah zdůvodni.
\tcblower
Rozvineme čitatel i příslušné faktoriály:
\[
 \binom nk=\frac{n!}{k!(n-k)!}
 =\frac nk\frac{(n-1)!}{(k-1)!(n-k)!}
 =\frac nk\binom{n-1}{k-1}.
\]
Dělení $k$ je dovoleno, protože $k\geq1$. Kombinatorická kontrola: komisi s jedním označeným vedoucím lze sestavit buď výběrem komise a poté vedoucího ($k\binom nk$), nebo výběrem vedoucího a pak ostatních členů ($n\binom{n-1}{k-1}$). Rovnost těchto počtů dává stejný vztah.
\end{prikladReseniBox}


\newpage
\blok{Samostatná práce a zásoba k procvičení}

Než začneš počítat, pojmenuj použitou vlastnost a ověř její indexy. C5 je nepovinné rozšíření: nestačí uvést polovinu celkového počtu bez vysvětlení, proč jsou oba počty stejné.

% Zdroj C1: import_hodiny/37_Vlastnosti_kombinacnich_cisel.tex, Procvičování, položka 1.
\begin{zadaniReseniBox}[title={Úloha C1}]
Zjednoduš \(\binom{12}{11}\).
\tcblower
Symetrií $\binom{12}{11}=\binom{12}{1}=12$.
\end{zadaniReseniBox}

% Zdroj C2: import_hodiny/37_Vlastnosti_kombinacnich_cisel.tex, Procvičování, položka 2.
\begin{zadaniReseniBox}[title={Úloha C2}]
Dopočti \(\binom73+\binom74\).
\tcblower
Horní indexy jsou stejné a spodní sousední. Pascalův vztah dává $\binom73+\binom74=\binom84=70$. Kontrola symetrií: $\binom73=\binom74=35$, tedy $35+35=70$.
\end{zadaniReseniBox}

% Zdroj C3: import_hodiny/37_Vlastnosti_kombinacnich_cisel.tex, Procvičování, položka 3.
\begin{zadaniReseniBox}[title={Úloha C3}]
Vyjádři \(\binom83+\binom82\) jediným členem.
\tcblower
Použijeme Pascalův vztah s $n=9,k=3$:
\[
\binom83+\binom82=\binom93=84.
\]
Kontrola: $56+28=84$.
\end{zadaniReseniBox}

% Zdroj C4: import_hodiny/37_Vlastnosti_kombinacnich_cisel.tex, Procvičování, položka 4.
\begin{zadaniReseniBox}[title={Úloha C4}]
Urči součet \(\binom60+\binom61+\cdots+\binom66\).
\tcblower
Jde o všechny velikosti podmnožin šestiprvkové množiny. Součet je $2^6=64$. Přímá kontrola: $1+6+15+20+15+6+1=64$.
\end{zadaniReseniBox}

% Zdroj C5: import_hodiny/37_Vlastnosti_kombinacnich_cisel.tex, Procvičování, položka 5.
\begin{zadaniReseniBox}[title={Úloha C5}]
Kolik podmnožin sedmiprvkové množiny má lichý počet prvků? Každý prvek lze vybrat nejvýše jednou.
\tcblower
Pevně zvolme jeden prvek P. Každé podmnožině změníme právě jeho přítomnost: není-li P vybrán, přidáme ho; je-li vybrán, odebereme ho. Změní se parita velikosti a opakování stejného kroku vrátí původní podmnožinu. Tím vytváříme dvojice jedné liché a jedné sudé podmnožiny. Všech podmnožin je $2^7=128$, lichých tedy $128/2=64$. Číselně také $\binom71+\binom73+\binom75+\binom77=7+35+21+1=64$.
\end{zadaniReseniBox}


\poznamka{Typické chyby: Pascalovo sčítání čísel s různými horními indexy; záměna symetrie za prohození $n$ a $k$; vynechání prázdné množiny v součtu. U C2 lze výsledek nezávisle spočítat symetrií jako $2\binom73$. Příště uspořádáme čísla do trojúhelníku a Pascalův vztah uvidíme i při počítání cest.}
\end{document}
